\documentclass[12pt,a4paper]{article}

% Encoding and fonts
\usepackage[utf8]{inputenc}
\usepackage[T1]{fontenc}
\usepackage{lmodern}

% Hyperlinks
\usepackage{hyperref}

% Listings for code
\usepackage{listings}
\usepackage{xcolor}

\lstdefinestyle{streem}{
  language=Ruby, % Streem syntax resembles Ruby
  backgroundcolor=\color{gray!10},
  basicstyle=\ttfamily\small,
  keywordstyle=\color{blue}\bfseries,
  commentstyle=\color{green!50!black}\itshape,
  stringstyle=\color{red!70!black},
  showstringspaces=false,
  frame=single,
  breaklines=true
}

% Tables
\usepackage{array}
\renewcommand{\arraystretch}{1.2}

% Title
\title{Parsing Streem grammar: a Shell-like programming language prototype}
\author{Maciej Matiaszowski}
\date{}

\begin{document}

\maketitle

\section{Screemer}

I speak \href{https://github.com/matz/streem}{Streem}, a prototype of a stream-based programming language.  
Screemer is a set of example programs in Streem, a beginning of something that should end as a Streem documentation project. 

\section{The Examples}

Run the examples with Docker to speak Streem:

\begin{lstlisting}[style=streem]
git clone https://github.com/Stagyrite/speakingstreem.git
cd speakingstreem
docker build --tag 'speakingstreem' .
\end{lstlisting}

\begin{itemize}
\item \href{fibonacci.strm}{Fibonacci numbers}
\item \href{languageVersions.strm}{Language versions (KVS demo)}
\end{itemize}

\section{Streem Language Documentation}

\subsection{Code comments}

All code comments are prefixed with a hash:

\begin{lstlisting}[style=streem]
# minimal program
\end{lstlisting}

\subsection{Types}

\begin{center}
\begin{tabular}{|c|c|}
\hline
\textbf{Value} & \textbf{Type} \\
\hline
"HELLO, world" & string \\
true, false & boolean \\
42 & integer \\
4.2 & float \\
{[}1,2,3{]} & array \\
x -> x & function \\
map, flatmap etc & built-in function \\
stdin, stdout etc & file descriptor \\
x | y & stream \\
\hline
\end{tabular}
\end{center}

\subsection{Assigning a variable}

\begin{lstlisting}[style=streem]
x = 42
s = "HELLO, world"
\end{lstlisting}

\subsection{If statements}

\begin{lstlisting}[style=streem]
# Output: Two
x = 2

if (x == 1) {
  s = "One"
} else if (x == 2) {
  s = "Two"
} else {
  s = "Unknown"
}

[s] | stdout
\end{lstlisting}

\subsection{Concatenate strings}

\begin{lstlisting}[style=streem]
# Output: HELLO, world

name = "world"
s = "HELLO, " + name
[s] | stdout
\end{lstlisting}

\subsection{Array}

\begin{lstlisting}[style=streem]
# Output:
# HELLO,
# world

x = ["HELLO,", "world"]
x | stdout
\end{lstlisting}

\section{Streem Library Documentation}

\subsection{puts()}

\begin{lstlisting}[style=streem]
# Output: HELLO, world
s = "HELLO, world"
puts(s)
\end{lstlisting}

\subsection{Stream example}

\begin{lstlisting}[style=streem]
x = ["HELLO,", "world"]
x | stdout
\end{lstlisting}

\subsection{Functions}

\begin{lstlisting}[style=streem]
# Output: HELLO, world
hello = { name ->
    "HELLO, " + name
}

s = hello("world")
print(s)
\end{lstlisting}

\subsection{Reverse an array}

\begin{lstlisting}[style=streem]
x = [1, 2, 3]
reverse(x) | stdout
# Output:
# 3
# 2
# 1
\end{lstlisting}

\subsection{Sequence of numbers}

\begin{lstlisting}[style=streem]
seq(3) | stdout
# Output:
# 1
# 2
# 3
\end{lstlisting}

\section{Math Library Documentation}

\subsection{Constants}

\begin{lstlisting}[style=streem]
print(PI)   # 3.141592...
print(E)    # 2.718281...
\end{lstlisting}

\subsection{Examples}

\begin{lstlisting}[style=streem]
print(ceil(45.54))   # 46
print(fabs(-2))      # 2
print(gcd(4,10))     # 2
print(trunc(9.13))   # 9
print(floor(2.7))    # 2
print(round(9.76))   # 10
print(pow(5,2))      # 25
print(sqrt(25))      # 5
\end{lstlisting}

\subsection{Trigonometry}

\begin{lstlisting}[style=streem]
print(hypot(3,4))    # 5
\end{lstlisting}

\end{document}
